% ============================================================================
%  CAPÍTULO II: HERRAMIENTAS UTILIZADAS Y METODOLOGÍA DE EVALUACIÓN
% ============================================================================

\section{Marco Teórico del Arsenal Tecnológico de Seguridad}

La evaluación de la postura de seguridad, la detección de vulnerabilidades y la verificación de las capacidades de respuesta defensiva demandan la integración coordinada de herramientas especializadas en las diferentes capas del modelo OSI \citep{nist_sp800_115}. A continuación se desarrollan los fundamentos conceptuales de las herramientas aplicadas en el presente ejercicio.

% -----------------------------------------------------------------------------
\subsection{Nmap (Network Mapper)}
% -----------------------------------------------------------------------------
\textbf{Nmap} constituye el estándar de facto en exploración de redes, auditoría de puertos y evaluación de superficies de ataque perimetral \citep{lyon_nmap_2009}. Opera mediante la emisión deliberada de paquetes IP crudos para determinar qué equipos están activos en el segmento, qué puertos TCP y UDP se encuentran abiertos, qué servicios y versiones de software responden, y qué sistema operativo subyace en el host.

Dentro de sus capacidades más avanzadas destaca el \textbf{Nmap Scripting Engine (NSE)}, el cual permite automatizar tareas complejas de auditoría mediante scripts desarrollados en lenguaje Lua. En este proyecto se utilizó la categoría \texttt{vuln} para interrogar la presencia de fallos conocidos y debilidades de configuración.

Asimismo, Nmap dispone de mecanismos para evaluación en entornos restringidos por cortafuegos e IPS, tales como la generación de señuelos aleatorios (\textit{Decoy Scanning} mediante \texttt{-D RND:20}), técnica que intercala la dirección IP real del atacante entre múltiples direcciones IP falsas generadas aleatoriamente, dificultando la correlación del origen de la sonda por parte de los sensores defensivos.

% -----------------------------------------------------------------------------
\subsection{Wireshark: Análisis Forense y de Protocolos de Red}
% -----------------------------------------------------------------------------
\textbf{Wireshark} es el analizador de protocolos de red y de paquetes de código abierto más utilizado a nivel mundial \citep{sanders_practical_2020}. Permite capturar e interactuar con el tráfico que atraviesa una interfaz de red física o virtual en tiempo real, operando en modo promiscuo.

Sus capacidades conceptuales clave comprenden:
\begin{itemize}
  \item \textbf{Inspección Profunda de Capas:} Decodificación completa de tramas desde la capa de enlace (Ethernet) hasta protocolos de aplicación (HTTP, DNS, SMTP, TLS, SMB, LDAP).
  \item \textbf{Filtros de Captura y Visualización:} Utilización de filtros BPF (\textit{Berkeley Packet Filters}) para aislar flujos sospechosos y reglas de visualización para reconstruir sesiones TCP (\textit{Follow TCP Stream}).
  \item \textbf{Verificación de Confidencialidad Criptográfica:} Comprobación empírica de si la información sensible (credenciales, cookies de sesión, mensajes de correo) transita en texto plano (HTTP, POP3, IMAP) o si se encuentra encapsulada dentro de canales criptográficos robustos (TLS 1.2/1.3).
\end{itemize}

% -----------------------------------------------------------------------------
\subsection{OWASP ZAP (Zed Attack Proxy)}
% -----------------------------------------------------------------------------
\textbf{OWASP ZAP} (Zed Attack Proxy) es una plataforma integral de pruebas de seguridad dinámicas para aplicaciones web (DAST -- \textit{Dynamic Application Security Testing}), mantenida por la comunidad global de OWASP \citep{owasp_top10_2021}. Funciona como un \textit{proxy} de interceptación intercalado entre el navegador del auditor y la aplicación web objetivo.

Entre sus módulos de operación destacan:
\begin{itemize}
  \item \textbf{Spider / Rastreador Web:} Mapea automáticamente la estructura de la aplicación, formularios y endpoints navegables.
  \item \textbf{Escaneo Pasivo:} Analiza las respuestas HTTP sin enviar cargas adicionales, identificando de manera no intrusiva problemas como la ausencia de cabeceras de seguridad (\texttt{Content-Security-Policy}, \texttt{X-Frame-Options}, \texttt{X-Content-Type-Options}) o cookies carentes de atributos \texttt{Secure} y \texttt{HttpOnly}.
  \item \textbf{Escaneo Activo:} Inyecta vectores de prueba dirigidos contra los parámetros de la aplicación para detectar inyecciones SQL (SQLi), Cross-Site Scripting (XSS), Local File Inclusion (LFI) y debilidades de control de acceso.
  \item \textbf{Clasificación de Alertas por Código de Color:} ZAP jerarquiza los hallazgos en cuatro niveles de riesgo: Rojo (\textit{High Risk}), Naranja (\textit{Medium Risk}), Amarillo (\textit{Low Risk}) y Azul (\textit{Informational}).
\end{itemize}

% -----------------------------------------------------------------------------
\subsection{Suricata IDS/IPS}
% -----------------------------------------------------------------------------
\textbf{Suricata} es un motor de seguridad de red de código abierto de última generación, capaz de operar como sistema de detección de intrusos (NIDS) y como sistema de prevención en línea (NIPS) \citep{oisf_suricata_2024}. Su arquitectura multi-hilo (\textit{multi-threaded}) optimiza el procesamiento paralelo del tráfico en interfaces de alta velocidad.

Suricata implementa inspección profunda de paquetes (DPI) y análisis contextual de protocolos (HTTP, TLS, DNS, SSH, SMB). Su motor compara las firmas del tráfico en tránsito contra reglas estandarizadas (como las reglas bancarias personalizadas o los conjuntos \textit{Emerging Threats}), pudiendo alertar pasivamente o descartar paquetes (\texttt{DROP}) en tiempo real cuando se acopla a las colas \texttt{NFQUEUE} del cortafuegos \texttt{nftables}.

% -----------------------------------------------------------------------------
\subsection{EveBox: Consola y Gestión de Telemetría en Tiempo Real}
% -----------------------------------------------------------------------------
\textbf{EveBox} es la interfaz web gráfica diseñada para la correlación, filtrado y gestión de eventos emitidos por el motor Suricata a través de su archivo de registro estructurado \texttt{eve.json} \citep{oisf_suricata_2024}. 

EveBox proporciona tableros en tiempo real con líneas temporales de incidentes, clasificación por severidad, gráficos de flujos de red y detalles minuciosos del paquete infractor (dirección IP origen y destino, payload decodificado, firma detonada y acción adoptada por el sensor).

% -----------------------------------------------------------------------------
\section{Evidencias de la Auditoría Dinámica Web y Monitoreo Defensivo}
% -----------------------------------------------------------------------------

En esta sección se incorporan las evidencias gráficas obtenidas durante las pruebas con las herramientas analizadas.

\subsection{Auditoría de Aplicaciones Web con OWASP ZAP}

Se ejecutó un escaneo dinámico en vivo mediante OWASP ZAP sobre la aplicación web independiente \textbf{Uyuni Way Web}, desplegada en el puerto \texttt{8085/tcp} del servidor objetivo.

La \textbf{FIGURA \ref{fig:escaneo-zaproxy}} muestra la ejecución del escaneo en vivo de ZAP sobre el puerto 8085, mientras que la \textbf{FIGURA \ref{fig:alerta-naranja-zap}} expone la matriz de alertas de severidad media (color naranja) descubiertas durante la auditoría.

\capturaAPA{fig:escaneo-zaproxy}{Escaneo Dinámico en Vivo con OWASP ZAP sobre Uyuni Way Web (Puerto 8085)}{escaneozaproxy.jpeg}{Captura de pantalla de la interfaz de OWASP ZAP ejecutando el escaneo activo contra el servicio en el puerto 8085.}

\capturaAPA{fig:alerta-naranja-zap}{Alertas de Nivel Medio (Color Naranja) Detectadas por OWASP ZAP}{alertanaranja.jpeg}{Reporte de vulnerabilidades de severidad media generadas por ZAP durante el escaneo sobre Uyuni Way Web.}

\subsection{Supervisión de Amenazas en Tiempo Real con EveBox}

La interacción de las herramientas ofensivas contra la infraestructura del host auditado fue capturada por el sensor Suricata y visualizada en la consola web EveBox.

La \textbf{FIGURA \ref{fig:evebox-alertas}} ilustra el panel de eventos de seguridad en tiempo real generado por EveBox en el puerto \texttt{5636/tcp}, confirmando la detección del tráfico anómalo generado durante las pruebas.

\capturaAPA{fig:evebox-alertas}{Panel de Eventos y Alertas de Seguridad en Tiempo Real en EveBox}{evebox.jpeg}{Consola gráfica de EveBox visualizando las alertas de intrusión y escaneo registradas por Suricata.}

\subsection{Comportamiento Defensivo del Servidor ante Escaneos Agresivos (Anti-Nmap)}

Durante las fases iniciales de auditoría, se ejecutó el comando estándar de escaneo de vulnerabilidades:
\begin{lstlisting}[style=terminal, caption={Comando Agresivo Inicial de Nmap}, label={lst:cmd-agresivo}]
sudo nmap -sV --script vuln 192.168.100.85
\end{lstlisting}

Dicho escaneo experimentó una retención prolongada estancándose en el \textbf{99.58\%} de ejecución debido a la acción defensiva del servidor web Nginx y del cortafuegos del host. Al detectar ráfagas agresivas de sockets lentos (asociadas a scripts DoS como Slowloris), las reglas del servidor aplicaron descartes silenciosos (\textit{silent drops}), obligando a Nmap a esperar indefinidamente por retransmisiones TCP.

La \textbf{FIGURA \ref{fig:antinmap}} evidencia este fenómeno de tarpitting defensivo registrado en la terminal del auditor.

\capturaAPA{fig:antinmap}{Estancamiento de Nmap NSE al 99.58\% por Efecto Tarpitting Defensivo en Nginx}{antinmap.jpeg}{Captura de pantalla de la terminal en Kali Linux evidenciando el bloqueo del escaneo agresivo de Nmap.}
